\chapter{Tâche 4 — Outillage du framework : extension VS Code (LSP et débogueur)}
\label{chap:t4}
% Entreprise : Free Mobile

% =====================================================================
%  Trame conseillée (réutilisée pour chaque tâche) :
%  1. Contexte & dette technique de départ
%  2. Objectif de la mission
%  3. Conception / choix techniques
%  4. Réalisation (extraits de code COMMENTÉS, schémas — PAS de code entier)
%  5. Difficultés rencontrées
%  6. Résultats & bénéfice business (réduction de la dette)
% =====================================================================

\section{Contexte et dette technique initiale}
Une fois le framework de test livré (tâche~3, chapitre~\ref{chap:t3}), un nouveau
frein est apparu : la rédaction des tests elle-même. Les scénarios sont décrits
en YAML, support sur lequel l'écriture ne bénéficie d'aucune assistance
contextuelle. Là où des langages comme Python, Lua ou C++ disposent d'un serveur
de langage (\gls{lsp}) offrant autocomplétion et analyse syntaxique au fil de la
frappe, le YAML n'est outillé que pour la validation de schémas. Il ne propose
pas d'assistance qui tienne compte du contexte — capable, par exemple, de savoir
si l'on est en train de saisir une instruction ou l'un de ses arguments.

Or le format de test n'est pas un simple document de configuration : c'est en
réalité un langage à part entière, bâti au-dessus du YAML. On y manipule des
\textbf{instructions} — équivalentes à des fonctions —, des \textbf{constantes}
ainsi que des \textbf{variables} qui constituent les arguments passés à ces
instructions. Faute d'outillage tenant compte de cette grammaire, les erreurs de
syntaxe n'étaient découvertes qu'à l'exécution et aucun moyen n'existait
d'observer pas à pas le déroulé d'un scénario. Cette friction freinait l'adoption
du framework par l'équipe : une dette « d'outillage ».

\section{Objectif de la mission}
La mission a consisté à développer une extension \emph{Visual Studio Code},
écrite en \textbf{TypeScript}, au service du framework de la tâche~3. Cette
extension réunit deux facettes complémentaires :
\begin{itemize}
  \item un \textbf{serveur de langage} (\gls{lsp}) apportant autocomplétion,
        diagnostics en temps réel et aide à la saisie. En s'appuyant sur la
        grammaire du format de test — instructions, constantes et variables —,
        il fournit une assistance \emph{contextuelle} : il sait distinguer
        l'endroit où l'on attend une instruction de celui où l'on saisit un
        argument, et propose les complétions correspondantes ;
  \item un \textbf{débogueur} permettant d'exécuter un test puis, lorsqu'une
        instruction échoue, d'inspecter l'état du \emph{Valo Engine} au moment
        précis de l'échec, ou de progresser \emph{pas à pas} pour comprendre le
        comportement interne du moteur en fonction des trames envoyées par le
        test.
\end{itemize}

Une fonctionnalité particulièrement utile en a découlé : la possibilité de lancer
le \emph{Valo Engine} en \textbf{mode débogage} et de progresser pas à pas
\emph{à l'intérieur même de son code C++}, tout en rejouant un scénario d'appel
(\emph{call flow}) étape par étape. Le diagnostic ne se limite ainsi plus au
test, mais descend jusqu'au code natif du logiciel sous test.

\section{Conception et choix techniques}
% TODO : un SCHÉMA d'architecture serait le bienvenu :
% \begin{figure}[H]\centering
%   \includegraphics[width=0.85\textwidth]{archi-extensions-vscode.png}
%   \caption{Architecture de l'extension et son dialogue avec le framework.}
%   \label{fig:archi-t2}
% \end{figure}
Le choix du \gls{lsp} répond à un principe de \textbf{découplage} entre l'éditeur
et le langage : la logique d'analyse et de complétion est portée par un serveur
indépendant, réutilisable au-delà de VS Code par tout éditeur compatible. Le
serveur exploite la grammaire du format de test — celle-là même qu'analyse le
parseur du framework (chapitre~\ref{chap:t3}) — pour produire des complétions et
des diagnostics cohérents avec ce que le framework acceptera réellement à
l'exécution.

Le débogueur, lui, repose sur le \textbf{Debug Adapter Protocol} (\gls{dap}),
protocole standardisé par VS Code pour piloter une session de débogage au moyen
d'échanges \emph{JSON-RPC}. Côté framework, son intégration s'appuie sur un
module qui contrôle le flot d'exécution du test — mise en pause, reprise,
progression pas à pas — tout en remontant à l'éditeur les données nécessaires au
diagnostic.

Les seules données internes au framework de test ne suffisent toutefois pas à
diagnostiquer l'ensemble des problèmes : il faut aussi connaître l'état interne
du moteur. Un \textbf{service} dédié — développé par David — a donc été mis en
place pour exposer l'état des variables internes du \emph{Valo Engine}.
L'extension dialogue avec ce service via \textbf{HTTP} et restitue ces valeurs
directement dans l'\gls{ide}, aux côtés de l'état du test.

\section{Réalisation}
% TODO : remplacer les extraits ci-dessous par des extraits réels du dépôt.
%        Ils sont fournis à titre ILLUSTRATIF (pas le code entier).
Côté \gls{lsp}, le serveur s'abonne aux requêtes de complétion émises par
l'éditeur. À partir de la position du curseur, il détermine le contexte
grammatical courant et renvoie la liste des propositions adaptées — instructions
disponibles ou variables déjà déclarées.

\begin{lstlisting}[language=Java,caption={Fourniture d'autocomplétion contextuelle via le LSP (extrait illustratif, TypeScript).},label={lst:t4-lsp}]
// Déclenché par l'éditeur à chaque demande de complétion
connection.onCompletion((params) => {
  // On déduit le contexte grammatical à partir de la position du curseur
  const contexte = analyserContexte(document, params.position);

  // Selon que l'on attend une instruction ou un argument,
  // on ne propose pas le même jeu de complétions
  return contexte === "instruction"
    ? instructionsSupportees()   // liste tirée du registry du framework
    : variablesDeclarees(document);
});
\end{lstlisting}

Côté débogueur, l'adaptateur \gls{dap} traduit les ordres de l'éditeur (poser un
point d'arrêt, exécuter l'instruction suivante) en commandes pilotant le flot
d'exécution du test, puis renvoie l'état courant — instruction en cours, état du
\emph{Valo Engine} obtenu via le service HTTP.

\begin{lstlisting}[language=Java,caption={Pas à pas : pose des points d'arrêt côté adaptateur de débogage (extrait illustratif, TypeScript).},label={lst:t4-dap}]
// Requête DAP : l'éditeur transmet les lignes où s'arrêter
protocol.onSetBreakpoints((args) => {
  // Les lignes proviennent du parseur du framework,
  // qui associe chaque instruction à sa position dans le YAML
  moteurDeTest.armerPointsArret(args.lines);

  // On confirme à l'éditeur les points d'arrêt effectivement posés
  return { breakpoints: args.lines.map(verifierLigne) };
});
\end{lstlisting}

\section{Difficultés rencontrées}
La première difficulté a tenu à la \textbf{documentation de l'API d'extension de
VS Code}. Si des exemples d'extensions existent, l'\gls{api} elle-même reste
inégalement documentée, et aucun standard ne s'impose réellement quant à la
manière de structurer le code. Les exemples disponibles sont fréquemment
\textbf{monolithiques}, ce qui offre peu de modèles d'architecture réutilisables
et a imposé un travail d'expérimentation pour aboutir à une organisation
maintenable.

La seconde difficulté a porté sur l'\textbf{intégration du \gls{dap}} au
framework de test. Il a fallu adapter le protocole — fondé sur des échanges
\emph{JSON-RPC} — à un module capable de contrôler le flot d'exécution d'un test
tout en remontant les données utiles au diagnostic. La restitution de l'état
interne du moteur, obtenue par un service HTTP distinct, a par ailleurs supposé
de coordonner deux canaux de communication (\gls{dap} et HTTP) au sein d'une même
session de débogage.

\section{Résultats et bénéfice business}
% TODO : quantifier si possible (temps d'écriture d'un test, erreurs évitées,
%        taux d'adoption). Relier au thème dette technique (chap.~\ref{chap:veille}).
\emph{[À rédiger : temps d'écriture d'un test réduit grâce à l'autocomplétion et
aux diagnostics en temps réel, erreurs de syntaxe détectées avant l'exécution,
diagnostic accéléré des comportements inattendus — jusque dans le code C++ du
moteur — et, in fine, meilleure adoption du framework par l'équipe. À relier au
thème de la dette technique, chapitre~\ref{chap:veille}.]}
